\answer{ex:20:1}{$t_{\text{w}}=\tau_r\ln\frac{1-L}{1-H}=16.8$ ساعة، $t_{\text{s}}=\tau_d\ln\frac HL=5.8$ ساعة، والدور $22.5$ ساعة؛ ويشدّ المولّدَ إلى $24$ ساعة التأرجحُ اليومي للعتبات، أي حلقة مقفلة الطور~\eqref{eq:pll}.}
\answer{ex:20:2}{(أ)~$L=0.111$؛ ومع $L=0.092$ كان النوم سيدوم $9.6$ ساعة. (ب)~بـ$22\%$ ($1/0.82=1.22$)، لا بـ$18\%$.}
\answer{ex:20:3}{$H=\frac{p-1}{p-1/q}=0.623$، $L=H/q=0.093$، حيث $p=e^{16/\tau_r}$، $q=e^{8/\tau_d}$. بعد الإيقاظ عقب $4$ ساعات ينزاح الطور إلى الأبد (النوم أبكر بـ$7.2$ ساعة)؛ ومع العتبات المتأرجحة يعود في يومين أو ثلاثة.}
\answer{ex:20:4}{(أ)~$r_\pm=\frac12\bigl(1\pm\sqrt{1-4k/w}\bigr)=0.947,\ 0.053$؛ $a_{\text{hi}}=3.51$، $a_{\text{lo}}=0.49$. (ب)~العمل $\approx0.33$ ثانية، والصمت $\approx0.59$ ثانية، والدور $\approx0.93$ ثانية، ونسبة العمل $0.36$.}
\answer{ex:20:5}{$a_{\text{hi}}/r_+<b<a_{\text{lo}}/r_-$: $3.71<b<9.20$. يخفض \nt{ACET} قيمة $b$ إلى ما دون $3.71$: فينتقل التقاطع إلى الفرع العلوي، وتستقر حالة العمل ($r\approx1$).}
\answer{ex:20:6}{$4.9$ و$6.8$ و$8.8$ و$5.5$ و$8.9$ ساعة عند $D=24$ و$32$ و$40$ و$48$ و$64$: المدة يحددها الطور اليومي لبدء النوم، لا الدَّين.}
\answer{ex:20:7}{بعد B يكون الخطأ على A في المتوسط $0.51$ (من $0.19$ إلى $1.08$ على $20$ مجموعة؛ وللاستجابة الصفرية $1$). ومع الخلط يكون الخطآن كلاهما أقل من $10^{-4}$.}
\answer{ex:20:8}{مسألة مفتوحة. التحجيم المنتظم يحفظ نسب الأوزان، لكنه لا يحفظ استجابات العصبون العتبي إلا إذا حُجّمت العتبة بالنسبة نفسها؛ أما الإعادات المخلوطة فتغيّر النسب. وثبات التماسات الكبيرة يناقض التحجيم المنتظم الخالص.}
